\chapter*{General Introduction}
\addcontentsline{toc}{chapter}{General Introduction}
Urban transportation has become increasingly problematic in modern cities, particularly in corporate settings where a large number of employees travel daily to the same workplace \cite{worldbank_transport}. The widespread use of private vehicles has led to congestion, higher commuting expenses, and a noticeable environmental impact due to increased carbon emissions \cite{iea_transport}. Within this context, ride-sharing emerges as a practical approach to reduce the number of vehicles in circulation while improving resource utilization \cite{agatz2012optimization}.

Despite its potential, current carpooling solutions often fall short in real-world usage. Many existing platforms depend on rigid matching strategies that require passengers to be located near a driver’s starting point. This constraint limits the number of possible matches and does not accurately reflect real commuting behavior, where individuals are spread across different areas and follow varying routes. Moreover, these systems rarely provide effective real-time coordination and generally do not leverage user activity data to improve matching outcomes.

These limitations become even more evident in enterprise environments. Although employees typically share a common destination, their schedules and departure locations differ significantly, making static matching approaches inefficient. Addressing this issue requires systems that can adapt to spatial distribution, support live interaction, and make use of historical data to improve decision-making.

This project was developed as part of a final-year internship conducted within a corporate environment, where the need for a more structured and efficient commuting solution was clearly observed. The objective was not limited to conceptual design but extended to building a fully operational platform capable of functioning in a real organizational setting.

To overcome the identified challenges, this work introduces a real-time carpooling platform designed specifically for enterprise use. The system relies on a centralized backend responsible for managing business logic and maintaining a consistent system state across all operations. The frontend acts primarily as a responsive interface that reflects backend updates without embedding independent decision logic.

One of the core aspects of the platform is its event-driven architecture, which allows continuous synchronization between users. By relying on WebSocket-based communication, the system propagates updates instantly, enabling participants to interact within a shared and dynamically updated environment. This approach eliminates the need for frequent polling while ensuring efficient and low-latency coordination.

In addition, an observability layer is integrated into the platform to monitor system behavior. Through administrative dashboards and system event tracking, this component provides visibility into platform operations, facilitating debugging, performance analysis, and validation in real deployment conditions.

From an intelligence perspective, the system combines rule-based mechanisms with machine learning techniques. The matching process is modeled as a route-based optimization problem, where the system evaluates the overlap between driver routes and passenger locations to determine feasible pickup points along the path. This method increases flexibility compared to traditional fixed-point matching while maintaining route consistency, and it is implemented using geospatial analysis techniques.

To complement this approach, a machine learning module analyzes historical commuting data to identify recurring travel patterns. Based on classification models, the system estimates the likelihood of user commutes at specific times and generates proactive ride suggestions. This predictive layer enhances usability by reducing manual interaction, while remaining independent from the core matching logic to preserve system reliability.

Together, these components form a hybrid intelligent architecture, where deterministic methods ensure control and correctness, and data-driven models introduce adaptability and personalization.

The proposed solution is developed as a full-stack application combining a Flask-based backend, a Flutter mobile interface, and a PostgreSQL database. It also relies on external geospatial services, including OpenStreetMap and OSRM, to handle route computation and location-related processing. Real-time interaction is achieved through WebSocket-based communication, allowing users to receive instant updates and remain continuously synchronized during system usage.

The primary goal of this project is to design and implement a scalable and intelligent carpooling platform that addresses the shortcomings of existing solutions. More specifically, the system is designed to:

\begin{itemize}
\item Enhance ride-matching flexibility by supporting dynamic pickup points along routes.
\item Enable continuous synchronization between users through an event-driven communication model.
\item Incorporate predictive features based on machine learning techniques.
\item Ensure consistency of system behavior through a centralized backend architecture.
\item Provide monitoring and observability tools to support validation and system analysis.
\end{itemize}

To achieve these objectives, the development followed an Agile methodology based on the Scrum framework. This approach allowed incremental design, implementation, and evaluation of system components throughout the project lifecycle.

The implementation process involved addressing several technical challenges, including system scalability, real-time responsiveness, data consistency, and seamless integration with external services. Particular emphasis was placed on maintainability and extensibility, ensuring that the platform can evolve to support future enhancements such as large-scale deployment and more advanced predictive models.

This project connects theoretical concepts with practical implementation by delivering a complete and operational system. By integrating real-time coordination, geospatial processing, and predictive analysis, the platform provides a foundation for more adaptive and efficient mobility solutions within enterprise environments.

The remainder of this report is organized as follows:

\begin{itemize}
\item \textbf{Chapter 1: Project Initiation} \\ 
Introduces the enterprise context, reviews existing solutions and their limitations, and presents the proposed approach along with the development methodology.

\item \textbf{Chapter 2: System Architecture and Design} \\ 
Describes the overall architecture, system layers, and key design decisions.

\item \textbf{Chapter 3: User Management and Authentication} \\ 
Details authentication mechanisms, user management, and administrative controls.

\item \textbf{Chapter 4: Ride Management and Real-Time System} \\ 
Explains ride management features and the integration of real-time communication.

\item \textbf{Chapter 5: Intelligent System (Matching and Machine Learning)} \\ 
Presents the route-based matching strategy and the machine learning model for commute prediction.

\item \textbf{Chapter 6: System Evaluation and Validation} \\ 
Evaluates the system using defined scenarios and analyzes its performance.
\end{itemize}

This work contributes to the field of intelligent transportation systems by introducing a hybrid approach that combines rule-based geospatial logic with predictive modeling. The results show that integrating deterministic methods with data-driven techniques can lead to practical and efficient mobility solutions suitable for real-world deployment.